\documentclass[11pt]{article}
\usepackage{icmlko}
\begin{document}

\title{씨앗보다 작은 손 --- 예보를 덜 바꾸는 변경이 자를 더 움직인다}
\author{월드모델 실험실}
\date{노트 476 · 2026년 8월 3일}
\maketitle

\begin{abstract}
노트 475 가 못 푼 것을 물었다 --- \emph{lr $0.16$ 은 학습 적합도를
$+0.0351$ 사는데 여섯 자 어디에서도 안 잃고 안 번다. 산 것이 어디로
갔나.} 가르는 법은 간단해 보였다: \textbf{유보 예보가 챔피언과 얼마나
다른가}. 거의 같으면 산 것이 학습 행에만 남은 것이다.

\begin{center}
\begin{tabular}{lrrr}
\hline
 & 학습 예보 & 유보 예보 & 학습$-$유보 \\
\hline
lr $0.16$ & $0.9935$ & $0.9907$ & $\mathbf{+0.0028}$ \\
\texttt{it} $400$ & $0.9922$ & $0.9917$ & $+0.0005$ \\
\textbf{씨앗만 바꿈}(챔피언끼리) & $0.9947$ & $\mathbf{0.9891}$ & $+0.0056$ \\
\hline
\end{tabular}
\end{center}

\textbf{등록한 문장은 반증됐다} --- \texttt{it} $400$ 의 유보 예보가
lr $0.16$ 보다 챔피언과 \emph{더} 비슷하다($0.9917 > 0.9907$).

\textbf{그런데 아래 줄이 더 크다.} 두 후보 모두 유보 예보가
\textbf{씨앗을 바꿀 때보다 덜 움직인다}($0.9907 \cdot 0.9917$ 대
$0.9891$). 그러면서 \texttt{it} $400$ 은 KR 에서 $-0.0085$ 를 일관되게
잃고 lr $0.16$ 은 아무 데서도 안 잃는다.
\end{abstract}

\begin{figure}[t]
\centering
\includegraphics[width=\textwidth]{figs/f1.pdf}
\caption{\textbf{가} --- 두 후보의 유보 예보 일치가 씨앗 잡음 바닥
(붉은 선)보다 위에 있다. \textbf{나} --- 덜 움직인 \texttt{it} $400$ 이
더 잃는다.}
\end{figure}

\section{크기가 아니라 일관성이다}

씨앗을 바꾸면 예보가 무작위로 흔들리고, 씨앗을 여럿 평균 내면 그 흔들림이
상쇄된다. \texttt{it} $400$ 이 움직이는 양은 그보다 작지만
\textbf{방향이 늘 같다} --- 그래서 씨앗을 짝지어 재면 상쇄되지 않고
$-0.0085$ 로 남는다.

\begin{quote}
\textbf{읽는 법} --- 변경의 중요도를 \textbf{예보가 얼마나 바뀌었나로
재면 안 된다.} 씨앗 잡음이 그보다 크다. \emph{같은 씨앗끼리 짝지어
재는 설계만이} 작고 일관된 손을 본다.
\end{quote}

이것이 이 사이클에서 씨앗 넷짜리 훑기가 세 번 틀린 이유이기도 하다
(노트 $451 \cdot 458 \cdot 473$). 신호가 \textbf{개별 예보의 잡음 바닥
아래}에 산다.

\section{그래도 lr 물음에는 반쯤 답한다}

``학습$-$유보'' 칸을 보면 방향이 맞다. \textbf{lr $0.16$ 의 불일치는
학습 쪽에 $5.6$배 더 쏠려 있다}($+0.0028$ 대 $+0.0005$). \texttt{it}
$400$ 은 학습과 유보를 \emph{거의 똑같이} 바꾼다 --- 함수 전체를 옮긴다.

\textbf{lr 이 산 적합도는 학습 행 쪽에 더 남는다.} 다만 그 쏠림이
씨앗끼리의 쏠림($+0.0056$)보다도 작아서, \emph{이 자로는 확정을 못
한다.} 노트 475 의 물음은 \textbf{반쯤 열린 채로 둔다}.

\section*{조항 --- 예보 변화량을 근거로 쓰지 않는다}

\begin{quote}
\textbf{조항} --- ``이 변경은 예보를 거의 안 바꾼다''를 \textbf{무해함의
근거로 쓰지 않는다.} 이 실험실의 실제 효과는 개별 예보의 씨앗 잡음보다
작다. 무해함은 \emph{짝지은 자}로만 보인다.
\end{quote}

\section*{남은 것}

\begin{itemize}
\item \textbf{lr 이 산 것의 행방}은 여전히 반쯤 열려 있다. 다음 자
후보는 \emph{나무 구조}다 --- 잎 수 · 쓰인 축 수 같은 것. 예보가 아니라
모형을 들여다보는 자다.
\item \textbf{씨앗을 더 쓰면 자가 세지나} --- 지금 짝 설계는 씨앗
여덟$\sim$열이다. 열여섯이면 $-0.0005$ 같은 것도 볼 수 있나.
\item 노트 475 의 건너뛰기 조항은 그대로 --- 이번 결과가 그것을 안
건드린다(적합도 이야기지 예보 이야기가 아니다).
\end{itemize}

\end{document}
